\section{The Note: One Vibe Experiencing Another}

A vibe does one thing: it experiences other vibes. That act of experiencing has a name, the note. This chapter is about the note, the only relationship in the entire theory, and about a small extra quality each note carries.

\subsection{Experiencing is the only relation}

In most pictures of the world there are many kinds of relationship: forces, fields, bonds, signals, causes. In this theory there is exactly one. A vibe experiences another vibe. That is the note. Every force, every bond, every signal you will ever meet in this walkthrough is built out of notes, vibes experiencing vibes, and nothing else.

When two cells of the crystal touch, the shared edge between them is a note. So the edges of the graph we built in Part II are notes. The crystal is a web of vibes (the cells) joined by notes (the edges). To experience and to be connected are the same thing here. Connection just is mutual experiencing.

\subsection{A note is not a third thing}

This is worth saying carefully, because it is easy to slip. The note is not a separate object sitting between two vibes, like a string tied between two balls. There is no string. The note is one vibe's act of experiencing another. It is a relationship, not a part. The whole theory has exactly one kind of stuff, the vibe, and the note is what vibes do, not an additional ingredient. When we draw an edge between two cells, the edge is a picture of an act, not a picture of a thing.

This keeps the theory honest about its monism. If notes were a second kind of stuff, we would have two ingredients, not one. They are not. There are vibes, and there is the experiencing they do. That is one substance and its activity.

\subsection{The fill: strength and sign of attention}

Each note carries a little extra structure, called its fill. The fill is the strength and sign of the attention along that note, how much, and in which direction, one vibe's tone counts toward the other. Like the tone, the fill is three-valued: it can pull in the same direction, pull in the opposite direction, or not pull at all.

Think of it as the character of a relationship. Some relationships are reinforcing: when my neighbor leans toward pleasure, that pushes me toward pleasure too. Some are opposing: my neighbor's pleasure pushes me toward pain. Some are indifferent: my neighbor's state does not move me. The fill records which kind each note is. It is the texture of the connection.

The fills are where the crystal's particular patterns get their shape. The same web of vibes with different fills behaves completely differently, just as the same network of people behaves differently depending on who reinforces and who opposes whom. When we later talk about memories, selves, forces, and fields, much of the specific content lives in the pattern of fills across the notes. The fills are the wiring.

\subsection{Will: the outward face}

One more small word, for completeness. We sometimes speak of a vibe's will. The will is just the vibe's tone regarded from the outside, what it shows to its neighbors, the state they read off it when they experience it. Tone is the feeling from the inside. Will is the same state as presented outward. They are not two things. They are one state seen from two sides, inner and outer.

We now have the full local picture: a vibe with a tone (its felt charge), experiencing its neighbors through notes (acts of experiencing), each note carrying a fill (the strength and sign of that attention), and each vibe presenting its tone outward as will. The only thing left is the rule that says how a vibe updates from all this. That is the next chapter, after we pause to make sure the vocabulary has not misled us.

\subsection{What sets the fills, and whether they change}

Where do the fills come from, and are they frozen? Fills are part of the structure, set as the crystal grows and its notes form. But they are not fixed forever. In living, learning systems they adapt with use: a connection that is repeatedly active strengthens, and one that is not weakens, the same principle that lets brains learn. This is how memory and habit and selves are laid down, as patterns written into the fills over time. So the fills are the malleable wiring of the world, given at first by the growth and then reshaped by experience, which is what makes learning, and a personal history, possible at all.

\textbf{Takeaway:} the note is the one and only relation, a vibe experiencing another vibe, drawn as an edge of the crystal. It is an act, not a third thing. Each note carries a three-valued fill, the strength and sign of the attention, which is the wiring that gives patterns their shape. Will is just the tone seen from outside.
